
% En esta sección se detallan los experimentos realizados para abordar las cuestiones propuestas en la sección \ref{Introduction}. En primer lugar, se introducen los detalles del diseño de la experimentación: la selección del modelo de segmentación, dataset elegido, procesos de entrenamiento y validación, etc. Después, se detalla el pipeline implementado para realizar la anotación semiautomática de las máscaras de ocupación, oclusión y área conducible. Finalmente, se aborda la metodología seguida para evaluar tanto la generación de las máscaras semánticas en \aclink{BEV}, como las detecciones 3D a partir de imágenes monoculares para la evaluación final del pipeline de anotación semiautomática.

This section details the experiments and implementations made to address the problems described in section \ref{Introduction}. On the first hand, the experimental design is introduced, tackling the model and dataset selection, how data augmentations are performed and the training and validation processes. Afterward, the implementation of the occupancy and occlusion masks preannotation pipeline is presented. Finally, the evaluation strategy of the pipeline is discussed to measure the quality of the resulting semantic masks and the monocular 3D detections for estimating the object's dimensions.

% Este proyecto puede dividirse en 3 bloques principales: experimentación con la segmentación semántica en BEV, diseño e implementación de un sistema de anotación semiautomático de máscaras de área conducible y un tercer bloque de experimentación para evaluar el rendimiento del sistema implementado.

Thus, this project can be divided into three main blocks: \aclink{BEV} semantic segmentation experimentation, design and implementation of the preannotation pipeline and a thrid block of how the system is evaluated.

\subsection{Segmentation experiment design: BEV2Seg\_2}
\label{bev2seg_2}

To address the hypotesis that "A semantic segmentation model trained with BEV images better segments planar elements" the process represented in Figure \ref{fig:beg2seg_2_flow} has been designed. It follows two main approaches: first, performing image segmentation on regular RGB images and then reprojecting them to \aclink{BEV}; second, applying \aclink{IPM} to the original images and then segmenting them with the model.

\begin{figure}[h!]
    \centering
    \includegraphics[width=\linewidth]{./images/metodology/bev2seg_2_flow.png}
    \caption{bev2seg\_2 diagram flow.}
    \label{fig:beg2seg_2_flow}
\end{figure}

Considering this flowchart, three main questions arises: (1) which segmentation model to use, (2) which dataset to train the models with, and (3) how to compare the train models and select the optimal output of the pipeline. 

\subsubsection{Segformer}
There are multiple techniques and strategies for tackling semantic segmentation in both regular and \aclink{BEV} images \ref{sota}.  

Several models could be chosen to address the proposed hypothesis. The current state of the art includes both \aclink{CNN} and \aclink{ViT}-based models that achieve competitive results. As shown in Figure \ref{tab:model_comparison}, there is no significant difference in accuracy and inference speed among the top-performing models. Moreover, many of these models have already been applied in the context of \aclink{ADS}. In this work, Segformer \cite{segformer} has been selected as the semantic segmentation model due to its balance between performance and efficiency. Additionally, it is integrated into the Huggingface \cite{huggingface} ecosystem, which provides an optimized, parallelized implementation, facilitating distributed training.  

\begin{table}[h]
    \centering
    \resizebox{\linewidth}{!}{%
    \begin{tabular}{l l c c c}
        \toprule
        \textbf{Model Name} & \textbf{Encoder} & \textbf{Params (M)} $\downarrow$ & \textbf{FPS} $\uparrow$ & \textbf{Cityscapes test mIoU ($\%$)} $\uparrow$ \\
        \midrule
        DDRNet-39       & -             & 32.3      & -         & 80.4 \\
        PIDNet-L        & -             & 36.9      & -         & 80.6 \\      
        DeeplabV3+      & ResNet-101    & 62.7      & 1.2       & 80.9 \\
        SETR            & ViT-Large     & 318.3     & 0.5       & 82.2 \\
        Segformer       & MiT-B4        & 64.1      & 3.0       & 83.8 \\
        \bottomrule
    \end{tabular}%
    }
    \caption{Comparison of different models. Results are obtained from \cite{DDRNet} \cite{PIDNet} \cite{segformer}.  }
    \label{tab:model_comparison}
\end{table}


The Segformer model consists on an hierarchical Transformer  encoder, which extract coarse and fine features, and a lightweight \aclink{MLP} decoder to directly fuse these multiscale features and predict the segmentation mask (Figure \ref{fig:segformer_architecture}). Segformer comes with a series of Mix Transformer encoders (MiT) that share the same architecture but have different sizes: from MiT-B0 as the lightweigtest encoder for realtime inference, to MiT-B5 for best performance.

\begin{figure}[h!]
    \centering
    \includegraphics[width=\linewidth]{images/metodology/segformer_architecture.png}
    \caption{Segformer architecture}
    \label{fig:segformer_architecture}
\end{figure}

The training strategy involves using pretrained encoders from ImageNet-1K, attaching an untrained segmentation head decoder, and fine-tuning the entire model for semantic segmentation of vehicular scenes. Since the objective is to train the models directly on \aclink{BEV} images, which differ significantly from standard perspective images, the encoder layers will remain trainable rather than being frozen.


\subsubsection{BEVDataset}
In order to train the model on the segmentation task, a valid dataset must be selected. There are several semantic segmentation datasets for the \aclink{ADS} context such as Cityscapes \cite{Cityscapes}, that defines $30$ semantic classes and provides $5.000$ frames with pixel-level high-quality annotations and $20.000$ weakly annotated images; KITTI \cite{KITTI}, which provices $400$ annotated images with a $0.5$ split for training and validation following the Cityscapes annotation format; ApolloScape \cite{ApolloScape}, that provides $146.997$ frames with corresponding pixel-level annonations and pose information for $25$ labels; or NuImages, a subset of NuScenes \cite{nuscenes}, that contains annotated images on $26$ different labels. For all of them, the ego pose and camera parameters metadata are provided. 

There are other benchmarks \cite{WildDash} \cite{CamVid} but in orther to train the models a rich dataset is needed. Despite Cityscapes and ApolloScape being good options for this task, NuImages has been selected as it is one of the most used datasets for \aclink{ADS} tasks, it provides very accurate ego poses and camera parameters, it has 3D annotations and it has a very good documentation.

NuImages contains around $93.000$ samples with aproximately $80\%$ reserved for the training set and $20\%$ for validation. Additionally, NuImages includes a private test set reserved for benchmark evaluations, whose annotations are not publicly available.

To train the models in the pipeline, a parser has been developed to convert NuImages into a sub-dataset named BEVDataset. This dataset has all front-camera images with NuImages annotations. Since the test annotations in NuImages are private, the validation set has been further splitted to ensure fair comparisons between models from different pipelines.

The conversion process is performed using a custom parser named "OLDatasets", which transforms NuImages samples into the structured ASAM OpenLABEL \footnote{\url{https://www.asam.net/standards/detail/openlabel/}} format, where metadata for each frame is stored. In the case of BEVDataset, images are reprojected into the \aclink{BEV} domain using the \aclink{VCD} library. This library provides tools to handle OpenLABEL annotations and manage both 2D and 3D data efficiently.

The "OLDatasets" parser extracts the camera parameters for each sample and computes a \aclink{LUT} to apply \aclink{IPM} reprojection. Using this data, semantic pixel masks are generated and reprojected along with the original images into the \aclink{BEV} space. Since this reprojection involves image warping, the interpolation method must be carefully chosen:
\begin{itemize}
    \item Linear interpolation is applied to images.
    \item Nearest neighbor interpolation is used for semantic masks to preserve pixel class integrity. 
\end{itemize}

The virtual \aclink{BEV} camera parameters are fixed: the \aclink{BEV} reprojection is generated using a regular grid with a cell spacing of 1 meter, covering a total distance of 30 meters in front of the camera and 1 meter behind it. The resulting images have a resolution of $1024 \times 1024$ pixels.

Finally, BEVDataset contains a total of $16.427$ images, distributed as shown in Figure \ref{fig:bev_dataset}.

\begin{figure}[h!]
    \centering
    \includegraphics[width=\linewidth]{images/metodology/BEVDataset.png}
    \caption{BEVDataset structure and mini set samples}
    \label{fig:bev_dataset}
\end{figure}


\subsubsection{Data augmentations}
Data augmentations in \aclink{BEV} are not an easy task...

\subsubsection{Validation and comparison}
Mean intersection over union (mIoU) or Jaccard index is used as the metric for evaluation and comparison between the two approaches.

\subsection{ Driveable area automatic annotation}
\label{aplication}

\subsubsection{Depth estimation}
\subsubsection{Scene PCD}
\subsubsection{Instance scene PCD}
\subsubsection{Instance BEV mask}

\subsection{Evaluation methodology}
\label{evaluacion}

\subsubsection{3D detections evaluations}
\subsubsection{BEV masks evaluation}
\hl{Groundtruth BEV masks could be generated from annotations}


